\clearpage
\fancyhead[L]{\small OBSERVATIONS ASTRONOMIQUES.}
\hypertarget{native-646-lower}{}
\begin{center}\small Göttingische gelehrte Anzeigen. Stück 129. Seiten 1281..1284. 14. August 1826.\\[3pt]\rule{.72\linewidth}{.3pt}\end{center}

\emph{Observations astronomiques faites à l'observatoire royal de Paris, publiées par le bureau des longitudes.} Paris 1825. Bei Bachelier. 402 Seiten im Folioformat.

Die Beobachtungen, die an der Pariser Sternwarte in den Jahren 1800..1809 angestellt wurden, sind in den Bänden der Connaissance des tems für 1808..1812 und 1823..1825 veröffentlicht worden. Die späteren sollen nach einem Beschluss des Bureau des longitudes gesondert gedruckt werden. Der vorliegende erste Band enthält, mit Ausnahme der Beobachtungen am Repetitionskreis von \textsc{Reichenbach} mit drei Fuß Durchmesser, diejenigen, die in den Jahren 1810..1819 von vier Beobachtern, den Herren \textsc{Bouvard}, \textsc{Arago}, \textsc{Mathieu} und \textsc{Nicollet}, angestellt wurden. Die für die Beobachtungen verwendeten Instrumente werden in der Einleitung beschrieben. Das von \textsc{Ramsden} begonnene und von \textsc{Berge} vollendete Mittagsfernrohr wurde im August 1803 aufgestellt. Es hat $7\frac{1}{2}$ Fuß Brennweite und 4 Zoll Öffnung. Die Achse ist 4 Fuß lang, die Vergrößerung nicht ganz hundertfach. Seitlich verschiebbares Okular und Beleuchtung durch die Achse, wie es jetzt allgemein üblich ist. Zur Justierung des Mittagsfernrohrs wurden zwei Meridianzeichen errichtet, das nördliche in 1364 Metern Entfernung auf dem Palais du Luxembourg, das südliche in 1840 Metern Entfernung an einer Pyramide in der Ebene von Montrouge. Als Zielpunkte dienen kreisrunde Löcher von $2\frac{1}{4}$ Zoll Durchmesser in waagerecht verschiebbaren Metallplatten. Das nördliche hebt sich gegen eine dahinter aufgestellte weiß gestrichene Eisenblechplatte ab, das südliche gegen den Himmel. Die scheinbaren Durchmesser betragen $9''07$ und $7''1$. Es gibt fünf senkrechte Fäden. Ihre Abstände gelten als genau gleich und werden von Sternen am Äquator in $17^{\rm s}36$ durchlaufen. Aus welchem Material die Fäden bestehen, wird nicht erwähnt. Die Beobachtungen an diesem Instrument nehmen mehr als die Hälfte des Bandes ein. Sie betreffen die Sonne, den Mond, die meisten Planeten und die Fixsterne von \textsc{Maskelyne}, selten andere Objekte. Von den Kometen von 1811 und 1819 sind einige untere Kulminationen verzeichnet. Was die Beobachtungen selbst betrifft, so genügen sie den Anforderungen, die die neuere Astronomie an selbstständige Beobachtungen mit Instrumenten von so außergewöhnlichen Abmessungen stellt, nicht ganz. Sie bieten keine ausreichenden Mittel, um zu untersuchen, mit welcher Genauigkeit die Meridianzeichen in der Meridianebene liegen. Bei der geringen Entfernung dieser Zeichen und der großen Brennweite des Fernrohrs müssen sie gegenüber den Fäden eine große Parallaxe zeigen und undeutlich erscheinen. Dazu kommt, dass sie insgesamt nicht oft zur Prüfung herangezogen wurden, etwa im Jahr 1810 zum ersten Mal am 11. Mai, nachdem in den früheren Monaten häufig vermerkt worden war, dass sie wegen des Nebels nicht zu sehen waren. Wenn die ungünstige Lage die Verwendung der Meridianzeichen so selten zuließ, wären häufige Beobachtungen von Zirkumpolarsternen doppelt notwendig gewesen. Doch schon der Polarstern wurde nur selten beobachtet und aufeinanderfolgende Kulminationen \SourceLeaf{647}oberhalb und unterhalb des Pols noch seltener. Zwar wurde oft geprüft, ob die Achse waagerecht liegt. Es fehlt jedoch gänzlich an Mitteln, um zu prüfen, ob die beiden Zapfen gleich dick sind, was man selbst bei Instrumenten der besten Instrumentenbauer nicht voraussetzen darf. Auch wurde nur selten geprüft, ob die optische Achse senkrecht zur Drehachse steht. In dieses Beobachtungstagebuch des Mittagsfernrohrs sind zahlreiche Beobachtungen von Sternbedeckungen eingefügt. Umso bedauerlicher ist es, dass die genannten Umstände bei der Bestimmung der absoluten Zeit stets eine geringe Unsicherheit zurücklassen, da die Astronomen gewohnt sind, die Angabe der geografischen Länge immer auf die Pariser Sternwarte zu beziehen.

\begingroup\rightskip=0pt plus 1em Der zweite Abschnitt des Werks besteht aus den Beobachtungen am $7\frac{1}{2}$ Fuß großen Mauerquadranten von \textsc{Bird}. Das Fernrohr ist ebenfalls $7\frac{1}{2}$ Fuß lang, hat $2\frac{1}{4}$ Zoll Öffnung und vergrößert 70..80-fach. Die beobachteten Objekte sind dieselben wie am Mittagsfernrohr. Man hat versucht, den Kollimationsfehler für eine Anzahl von Punkten aus zahlreichen Beobachtungen von Fixsternen zu bestimmen, wobei man für deren Deklinationen den Mittelwert aus den Angaben mehrerer Astronomen zugrunde legte. Es erübrigt sich, den Rang anzudeuten, den die Beobachtungen mit diesem Instrument beim gegenwärtigen Stand der praktischen Astronomie haben. Dieser Quadrant dient für die südliche Hälfte des Meridians. An der Nordseite ist ein zweiter, fünf Fuß großer Quadrant von \textsc{Sisson} aufgehängt, derselbe, den einst \textsc{Lalande} in Berlin benutzte. Der vorliegende Band enthält jedoch keine Beobachtungen mit diesem Instrument. --- Seit 1823 ist ein Mauerkreis von \textsc{Fortin} in Gebrauch, dessen Beschreibung wir im nächsten Band erwarten dürfen.\par\endgroup

Der dritte Abschnitt enthält die Beobachtungen mit der parallaktischen Maschine von \textsc{Bellet}, die in einem abgetrennten Teil der Sternwarte aufgestellt ist. Die Durchmesser ihrer beiden Kreise betragen 13 Zoll. Das Fernrohr ist drei Fuß lang, hat beinahe $2\frac{1}{2}$ Zoll Brennweite und vergrößert 40..50-fach. Die Beobachtungen betreffen teils die im Zeitraum von 1810..1819 erschienenen Kometen, teils den Mondfleck Manilius, der zur Untersuchung der Libration des Mondes beobachtet wurde. Die sehr wertvollen Ergebnisse dieser Untersuchung wurden bekanntlich schon vor mehreren Jahren in der Connaissance des tems veröffentlicht. Dieses Instrument wurde 1823 an die Sternwarte in Marseille abgegeben und durch ein großes Äquatoreal von \textsc{Gambey} ersetzt, dessen Beschreibung im nächsten Band folgen soll.
